% =====================================================================
% Appendix: construction and validation of difficulty_v1.
% Every number recomputed from the frozen response matrix by
% difficulty_appendix_stats.py, not copied from an earlier note.
% =====================================================================

\subsection{Construction of \texttt{difficulty\_v1}}
\label{app:difficulty}

\paragraph{Motivation.}
The examinations publish a level ordering, but that ordering is a curriculum sequence. It says
which material a candidate meets first, not which items are hard for a reasoning system, and the
two turn out to be nearly independent: Cram\'er's $V$ between the empirical bands defined below
and the published level labels is $0.174$. We therefore derive difficulty from behaviour.

\paragraph{Response matrix.}
Seventeen frozen systems answer all $10{,}198$ items once each under the protocol of
Appendix~\ref{app:hardware}, at temperature $0$ with seed $202607$. Responses from which no
option letter can be recovered count as incorrect, so every cell of the matrix is defined. The
matrix is frozen and its SHA-256 is
\texttt{bb093a59d00162f9d50cc67da4f0ee0748e326325badb2bb1747d90c4656e2be}. All bands in this
paper are derived from that exact artefact.

\paragraph{Why the average is weighted rather than flat.}
A flat mean over seventeen systems is not a neutral choice, because the panel is not balanced by
construction: ten systems are proprietary and seven have open weights, and the open-weight group
is on average substantially weaker. Under a flat mean the proprietary majority decides band membership, and an
item that every proprietary system solves is labelled easy even when no open-weight system does.
We therefore partition the panel by access, average within each group, and weight the two
group means equally. Writing $\mathcal{G}_{\mathrm{prop}}$ for the ten proprietary systems and
$\mathcal{G}_{\mathrm{open}}$ for the seven open-weight ones, and $c_{m,i}\in\{0,1\}$ for whether system $m$ answers item
$i$ correctly, the consensus pass rate is
\begin{equation}
s_i \;=\; \frac{1}{2}\sum_{k=1}^{2}\frac{1}{|\mathcal{G}_k|}\sum_{m\in\mathcal{G}_k} c_{m,i},
\label{eq:consensus}
\end{equation}
and the band is
\begin{equation}
\mathrm{band}(i)=
\begin{cases}
\textsc{easy}   & s_i \geq 2/3,\\
\textsc{hard}   & s_i \leq 1/3,\\
\textsc{medium} & \text{otherwise.}
\end{cases}
\label{eq:band}
\end{equation}
This yields $6{,}578$ easy, $2{,}183$ medium and $1{,}437$ hard items. The thresholds are the
natural thirds of the unit interval and were fixed before any band-conditioned result was
computed.

\paragraph{Alternative rules.}
Equation~\ref{eq:consensus} is one of three groupings we computed, and
Table~\ref{tab:difficulty-rules} reports all of them. Rule~A takes a flat mean over all
seventeen systems. Rule~B splits the panel three ways, into proprietary, open-weight reasoning
and finance-specialized, and weights the three equally. Rule~C is
Equation~\ref{eq:consensus} and is the one used throughout the paper.

The three rules agree closely. Rule~C and Rule~A place $95.8$ percent of items in the same band
and their consensus scores correlate at Spearman $0.998$. Rule~C and Rule~B agree on $90.8$
percent at Spearman $0.991$. Rule~B is the most aggressive because giving the five
finance-specialized systems a full third of the weight pushes items they fail into the hard band,
which inflates that band to $1{,}727$ items. We prefer Rule~C because the two-way split follows the licence boundary, which is a fact about
the system rather than a judgment of ours, whereas the three-way split additionally asserts that
finance-specialized systems form a category deserving equal weight with the other two.

\begin{table}[t]
\centering\small
\begin{tabular}{llrrr}
\toprule
\textbf{Rule} & \textbf{Grouping} & \textbf{Easy} & \textbf{Medium} & \textbf{Hard} \\
\midrule
A & Flat mean over 17          & 6{,}826 & 2{,}025 & 1{,}347 \\
B & Three groups, equal weight & 6{,}757 & 1{,}714 & 1{,}727 \\
\textbf{C} & \textbf{Two groups, equal weight} & \textbf{6{,}578} & \textbf{2{,}183} & \textbf{1{,}437} \\
\bottomrule
\end{tabular}
\caption{Band sizes under the three candidate rules. Rule~C is \texttt{difficulty\_v1} and is
used throughout the paper. Rules A and B are reported here as robustness checks and their labels
ship with the dataset.}
\label{tab:difficulty-rules}
\end{table}

\paragraph{Threshold sensitivity.}
Table~\ref{tab:difficulty-thresholds} varies the two cut points around the chosen thirds. Moving
them by $\pm 0.03$ leaves $96.8$ to $97.2$ percent of items in the same band. Widening to
$0.75/0.25$ or narrowing to $0.60/0.40$ still preserves about $90$ percent. The partition is
therefore not an artefact of the particular cut points. It is nonetheless a partition of a
continuous quantity, and $609$ items, $6.0$ percent, sit within $0.02$ of a cut point, so
per-item band membership near the boundary should not be treated as categorical. Analyses in the
main text that depend on the boundary are reported against the continuous score as well.

\begin{table}[t]
\centering\small
\begin{tabular}{lrrrr}
\toprule
\textbf{Cut points} & \textbf{Easy} & \textbf{Medium} & \textbf{Hard} & \textbf{Same band} \\
\midrule
$0.67/0.33$ (used) & 6{,}578 & 2{,}183 & 1{,}437 & --- \\
$0.70/0.30$        & 6{,}385 & 2{,}467 & 1{,}346 & $97.2\%$ \\
$0.65/0.35$        & 6{,}796 & 1{,}858 & 1{,}544 & $96.8\%$ \\
$0.60/0.40$        & 7{,}220 & 1{,}218 & 1{,}760 & $90.5\%$ \\
$0.75/0.25$        & 5{,}830 & 3{,}214 & 1{,}154 & $89.9\%$ \\
\bottomrule
\end{tabular}
\caption{Sensitivity of \texttt{difficulty\_v1} to the two cut points, holding
Equation~\ref{eq:consensus} fixed. The final column is the share of items that keep the band they
receive under the cut points actually used.}
\label{tab:difficulty-thresholds}
\end{table}

\paragraph{Reliability.}
Split-half reliability is estimated by repeatedly partitioning the seventeen systems at random
into halves of eight and nine, scoring every item under each half, and correlating the two
resulting orderings. Over $200$ random splits the mean Spearman correlation is $0.821$ with a
standard deviation of $0.024$ and a range of $0.737$ to $0.849$, giving a Spearman-Brown
corrected reliability of $0.902$. Removing any single system and rebuilding the score leaves the
ordering essentially unchanged: the worst leave-one-out correlation with the full-panel score is
$0.983$, obtained by dropping Hawkish-8B. No single grader determines the partition.

\paragraph{Relation to the published labels.}
Cram\'er's $V$ between the bands and the exams' own level labels is $0.174$, and between the
bands and the examination identity, CFA against FRM, it is $0.056$. The empirical axis is
therefore neither a restatement of the curriculum sequence nor a proxy for which examination an
item came from. Table~\ref{tab:difficulty-by-stage} gives the band composition of each stage.
CFA Level~II is the hardest stage on this axis with a quarter of its items in the hard band, and
FRM Part~I and CFA Level~I are the easiest, which is not the ordering the published level labels
would predict, since Level~III sits above Level~II in the curriculum but below it here.

\begin{table}[t]
\centering\small
\begin{tabular}{lrrrr}
\toprule
\textbf{Stage} & \textbf{Items} & \textbf{Easy} & \textbf{Medium} & \textbf{Hard} \\
\midrule
CFA Level I   & 4{,}763 & $72.9\%$ & $17.3\%$ & $9.8\%$  \\
CFA Level II  & 2{,}215 & $46.5\%$ & $28.0\%$ & $25.5\%$ \\
CFA Level III & 653     & $57.1\%$ & $23.0\%$ & $19.9\%$ \\
FRM Part I    & 1{,}652 & $71.7\%$ & $19.6\%$ & $8.7\%$  \\
FRM Part II   & 915     & $56.4\%$ & $29.1\%$ & $14.5\%$ \\
\bottomrule
\end{tabular}
\caption{Band composition by examination stage under \texttt{difficulty\_v1}.}
\label{tab:difficulty-by-stage}
\end{table}

\paragraph{Circularity control.}
Because the bands come from model consensus, scoring a system against a partition it helped
define would flatter it. Every band-conditioned result in the main text is therefore computed
leave-one-out: for system $m$, the bands are rebuilt from Equation~\ref{eq:consensus} over the
other sixteen systems, and $m$ is scored against that partition. This is why the per-model hard
band sizes in Table~\ref{tab:leaderboard} differ slightly from the $1{,}437$ reported here.
Appendix~\ref{app:circularity} repeats the analysis with DeepSeek-R1 and GPT-4o removed from the
panel entirely, since those two systems also serve as the backbones of the retrieval study in
Section~\ref{sec:analysis}.

\paragraph{Distribution of the consensus score.}
The score is strongly left-skewed, with a median of $0.807$ and a tenth percentile of $0.221$.
Ninety percent of items reach a consensus pass rate of $1.0$ or below by construction, and the
upper quartile begins at $0.950$, so the easy band is genuinely easy for this panel rather than
merely above a threshold. The hard band is the thin left tail, and its accuracy behaviour is the
subject of RQ1.

\subsection{Statistical Procedures}
\label{app:tests}

Every comparison in the paper is paired at the item level, because all systems answer the same
$10{,}198$ items. We list the tests used and the null each assumes.

\paragraph{Accuracy against chance.}
CFA items carry three options and FRM items four, so a system's hard-band score under random
responding is a sum of Bernoulli trials with differing success probabilities, that is a
Poisson-binomial rather than a binomial. We evaluate its distribution exactly by dynamic
programming over the per-item probabilities, which is tractable at the few hundred items per
band, and report the one-sided tail. Across the seventeen systems we control the false discovery
rate with Benjamini-Hochberg at $q<0.05$. Reporting a bare count of systems below chance would
not distinguish a real deficit from sampling noise, which matters because on the answerable
subset each system's hard band holds only about $390$ items.

\paragraph{Error concentration.}
The natural null is that a system's errors fall uniformly across the distractors of the item it
missed. A single constant floor is not valid, because the number of distractors varies with the
examination and the number of erring systems varies with the item, and the expected maximum of a
uniform allocation is above the mean at small counts. For each item we therefore enumerate all
allocations of its observed error count across its own distractors and take the expected modal
share as that item's null. The paired difference between observed and expected concentration is
tested with a Wilcoxon signed-rank test, which assumes only symmetry of the paired differences
and not normality.

\paragraph{Paired accuracy comparisons.}
Retrieval and verification conditions in Section~\ref{sec:analysis} are compared with an exact McNemar
test on the discordant pairs, with Benjamini-Hochberg correction across the conditions within each
family. McNemar is the appropriate test here because the conditions share items and differ only in the
evidence supplied, so the discordant counts carry all the information.

\paragraph{Association and reliability.}
Association between the empirical bands and the categorical labels published with the
examinations is measured with Cram\'er's $V$. Rank agreement between two scorings of the same
items, used for split-half reliability, leave-one-out stability and public against held-out
ranking, is measured with Spearman's $\rho$, and split-half values are corrected with the
Spearman-Brown formula. Proportions from the audits carry Wilson score intervals rather than
normal approximations, which behave badly at the small counts and extreme proportions that
appear in the calibration set.

\paragraph{What we do not claim.}
The seventeen systems are not a random sample from a population of language models, so intervals
computed across systems describe this panel and not models in general. Items are treated as the
sampling unit throughout, and where an item belongs to a multi-question case vignette its
siblings are not independent, which the tests above do not model. This affects the effective
sample size of item-level tests in the direction of overstating significance, and is the reason
we report effect sizes such as the paired median excess alongside every $p$-value rather than
relying on the $p$-value alone.

\paragraph{Released artefacts.}
Each record carries \texttt{difficulty}, \texttt{consensus\_pass\_rate},
\texttt{difficulty\_version}, and the labels under Rules A and B as
\texttt{difficulty\_band\_flat\_17} and \texttt{difficulty\_band\_3group}. The panel membership,
the per-model output hashes and the response-matrix hash ship alongside, so the bands can be
recomputed independently or replaced with a different rule without rerunning any model.
